% !TEX program = xelatex
% (c) Кафедра системного программирования, 2026

\documentclass[a4paper,11pt]{article}

\usepackage[
    a4paper,
    top=1.25cm,
    bottom=1.20cm,
    left=1.75cm,
    right=1.45cm
]{geometry}

\usepackage{fontspec}
\usepackage{polyglossia}
\setdefaultlanguage[babelshorthands=true]{russian}
\setotherlanguage{english}

\setmainfont{FreeSerif}
\setmonofont{DejaVuSansMono}

\usepackage{microtype}
\usepackage{parskip}
\setlength{\parindent}{0pt}
\setlength{\parskip}{0.12em}

\usepackage{titlesec}
\titlespacing*{\section}{0pt}{0.48em}{0.20em}
\titlespacing*{\subsection}{0pt}{0.32em}{0.12em}

\usepackage{xurl}
\usepackage{hyperref}
\usepackage{bookmark}
\usepackage{array}
\usepackage{booktabs}
\usepackage{float}

\hypersetup{
    colorlinks=true,
    linkcolor=black,
    urlcolor=blue,
    citecolor=black,
    breaklinks=true
}

\urlstyle{same}

\usepackage{titling}
\pretitle{\begin{center}\bfseries\Large}
\posttitle{\par\end{center}}
\preauthor{\begin{center}\normalsize}
\postauthor{\par\end{center}}
\setlength{\droptitle}{-0.9cm}

\renewcommand{\refname}{Список использованных источников}

\begin{document}

\raggedbottom

\title{Анализ графов}
\author{Пентегова Мария Сергеевна}
\date{}
\maketitle

\begin{flushright}
Группа: \emph{25.Б41-мм} \\
Кафедра: \emph{системного программирования} \\
Научный руководитель: \emph{Григорьев Семён Вячеславович} \\
Номер семестра практики: \emph{2}
\end{flushright}

\section{Введение}

GraphBLAS позволяет формулировать алгоритмы анализа графов в терминах
операций линейной алгебры над разреженными матрицами и векторами.
Цель данной работы --- реализовать обход графа в ширину (BFS) в
классическом и матрично-алгебраическом вариантах и исследовать
зависимость их производительности от структуры входного графа.

\section{Постановка задачи и реализация}

Реализованы четыре варианта BFS: классический BFS с построением
массива родителей (parent-BFS), классический multisource BFS,
GraphBLAS BFS с вычислением уровней и GraphBLAS multisource BFS.
GraphBLAS-варианты реализованы с использованием
SuiteSparse:GraphBLAS.

В классическом варианте граф хранится в формате CSR (Compressed
Sparse Row). Для вершины $v$ её соседи находятся в диапазоне
\[
[\texttt{row\_ptr}[v],\texttt{row\_ptr}[v+1])
\]
массива \texttt{col\_idx}. BFS использует очередь, а при первом
посещении вершины $u$ записывается
$\texttt{parent}[u]=v$. В multisource-варианте очередь изначально
содержит несколько исходных вершин. Сложность классического BFS
составляет $O(|V|+|E|)$.

Классическая реализация написана непосредственно на CSR, а не взята
из сторонней библиотеки. Это связано с организацией эксперимента:
один и тот же CSR используется как вход для классического BFS и как
исходное представление при построении матрицы GraphBLAS. Благодаря
этому обе реализации работают с одним и тем же набором рёбер и не
требуют дополнительного преобразования графа в формат сторонней
библиотеки.

Такой вариант позволяет отделить сравнение способов выполнения BFS
от различий библиотечных реализаций. Готовые реализации BGL или
LAGraph могут использовать собственные структуры данных и
оптимизации хранения графа. При их непосредственном использовании
эксперимент сравнивал бы не только классический и
матрично-алгебраический подходы, но и конкретные библиотечные
реализации. Поэтому в данной работе классический BFS сделан
намеренно простым: используется обычная очередь и последовательный
просмотр CSR. Его результаты одновременно используются для проверки
корректности GraphBLAS-вариантов.

В GraphBLAS граф представляется булевой матрицей смежности $A$, а
текущий фронт --- разреженным булевым вектором $f_k$. Следующий фронт
вычисляется матрично-векторной операцией
\[
f_{k+1}=f_kA
\]
на булевом semiring \texttt{LOR-LAND}. Уже посещённые вершины
исключаются с помощью маски. Номер итерации определяет уровень
вершины в BFS. В multisource-варианте начальный вектор содержит
несколько источников.

Таким образом, классическая реализация представляет собой
контрольный вариант последовательного BFS, а GraphBLAS-реализация
выражает тот же обход через операции над разреженными матрицами и
векторами. Такой подход соответствует назначению GraphBLAS
\cite{graphblas-api}.

Программа написана на C11. Для измерения времени используется
\texttt{clock\_gettime(CLOCK\_MONOTONIC)}. Перед серией измерений
GraphBLAS-вариантов выполняется прогревочный запуск. Каждый вариант
запускается 10 раз; для каждой серии вычисляются среднее, минимальное
и максимальное время.

Матрица GraphBLAS импортируется из CSR один раз до начала измерений.
Время её построения выводится отдельно и не включается во время BFS.
Таким образом, измеряется выполнение алгоритма на уже подготовленном
представлении графа, а стоимость однократного преобразования
учитывается отдельно.

\section{Эксперимент}

Для экспериментов выбраны семь графов из SuiteSparse Matrix Collection,
содержащей набор разреженных матриц для тестирования алгоритмов
работы с разреженными данными \cite{ssmc}.

Для однократного BFS источником выбирается вершина максимальной
степени. Такой выбор обеспечивает воспроизводимое определение
источника для каждого графа. Для multisource BFS используются четыре
источника с индексами
$0$, $n/4$, $n/2$ и $3n/4$.

\begin{table}[H]
\centering
\small
\setlength{\tabcolsep}{3.5pt}
\begin{tabular}{lrrrr}
\toprule
Граф & $|V|$ & $|E|$ & Тип & $\deg(s)$ \\
\midrule
com-Orkut        & 3.07M  & 117M   & социальный  & 33313 \\
soc-LiveJournal1 & 4.85M  & 69M    & социальный  & 38660 \\
road\_usa        & 23.95M & 28.85M & дорожный    & 4 \\
roadNet-CA       & 1.97M  & 2.77M  & дорожный    & 3 \\
web-Google       & 0.92M  & 5.10M  & веб-граф    & 6332 \\
cit-Patents      & 3.77M  & 16.52M & цитирование & 852 \\
com-Amazon       & 0.33M  & 0.93M  & социальный  & 546 \\
\bottomrule
\end{tabular}
\caption{Характеристики исследуемых графов.}
\label{tab:graphs}
\end{table}

Перед измерениями для GraphBLAS-вариантов выполняется один
прогревочный запуск. Затем каждый из четырёх вариантов BFS
запускается 10 раз. Время измеряется в микросекундах; для каждой
реализации вычисляются среднее, минимальное и максимальное значения.

\begin{table}[H]
\centering
\small
\setlength{\tabcolsep}{2.7pt}
\begin{tabular}{lrrrr}
\toprule
Граф & CP & CM & GB-L & GB-M \\
\midrule
com-Orkut        & 2652628 & 2700768 & 486500   & 420307 \\
soc-LiveJournal1 & 1915627 & 1976468 & 358255   & 331458 \\
road\_usa        & 1451387 & 1588466 & 25839879 & 22121848 \\
roadNet-CA       & 127543  & 110558  & 663630   & 670218 \\
web-Google       & 100233  & 80683   & 101149   & 106972 \\
cit-Patents      & 17500   & 4730    & 181639   & 46757 \\
com-Amazon       & 38967   & 31974   & 61923    & 62144 \\
\bottomrule
\end{tabular}
\caption{Среднее время BFS по 10 запускам, мкс.
CP --- Classic Parent, CM --- Classic Multisource,
GB-L --- GraphBLAS Level, GB-M --- GraphBLAS Multisource.}
\label{tab:results}
\end{table}

\section{Анализ результатов}

На \texttt{com-Orkut} и \texttt{soc-LiveJournal1} GraphBLAS Level
показал меньшее среднее время: 486\,500 и 358\,255 мкс против
2\,652\,628 и 1\,915\,627 мкс для CP. Отношение времени CP к GB-L
составляет примерно 5.45 и 5.35 соответственно.

На \texttt{road\_usa} и \texttt{roadNet-CA}, напротив, GraphBLAS Level
занял 25\,839\,879 и 663\,630 мкс против 1\,451\,387 и
127\,543 мкс для CP. На \texttt{web-Google} времена близки:
100\,233 и 101\,149 мкс. На \texttt{cit-Patents} и
\texttt{com-Amazon} классический вариант также оказался быстрее.

Таким образом, в исследованной выборке производительность зависит от
структуры графа. Одним из возможных факторов является ширина
BFS-фронтов: широкие фронты позволяют выполнить большой объём работы
матричной операцией, тогда как при узких фронтах накладные расходы
GraphBLAS могут становиться существеннее. В текущем эксперименте
размеры фронтов отдельно не измерялись, поэтому это объяснение
является гипотезой для дальнейшей проверки.

Полученные результаты также показывают, что само использование
матрично-алгебраического представления не гарантирует уменьшения
времени на каждом типе графа. Эффект зависит от того, насколько
структура конкретного графа и возникающие в BFS фронты подходят для
обработки разреженными матричными операциями.

\section{Тестирование}

Корректность реализации проверяется с помощью CTest и
\texttt{assert.h}. Тесты выполняются на небольших графах и проверяют
посещение вершин, вычисление уровней, построение родительского
дерева, multisource BFS и соответствие результатов классической и
GraphBLAS-реализаций.

Классический parent-BFS используется как контроль при проверке
GraphBLAS-результатов: для одних и тех же источников сравниваются
достижимость вершин и соответствующие расстояния. Все тесты проходят
успешно: \texttt{100\% passed}.

\section{Заключение}

В работе реализованы четыре варианта BFS на C: два классических
варианта на CSR и два варианта на SuiteSparse:GraphBLAS. Классическая
реализация выполнена самостоятельно как простая контрольная версия
на том же CSR, из которого строится GraphBLAS-матрица. Это позволяет
проводить сравнение на общем представлении входного графа и не
включать в эксперимент дополнительные различия между внутренними
форматами сторонних библиотек.

Эксперименты проведены на семи графах из SuiteSparse Matrix Collection.
Для каждого графа выполнены прогревочный запуск и десять измерений
каждого варианта. Построение GraphBLAS-матрицы выполняется один раз
и измеряется отдельно.

Результаты показывают зависимость производительности от структуры
графа. На \texttt{com-Orkut} и \texttt{soc-LiveJournal1} GraphBLAS
Level выполнялся примерно в 5.45 и 5.35 раза быстрее CP, тогда как
на \texttt{road\_usa}, \texttt{roadNet-CA}, \texttt{cit-Patents} и
\texttt{com-Amazon} меньшее время наблюдалось у классического
варианта. На \texttt{web-Google} результаты были близки.

Для дальнейшего эксперимента целесообразно измерять размер фронта на
каждом уровне BFS, число обработанных рёбер и стоимость отдельных
GraphBLAS-операций. Это позволит проверить связь структуры фронтов с
наблюдаемой разницей производительности.

Репозиторий проекта:

\url{https://github.com/MariaPentegova/graph_blas_bfs-2}

\begin{thebibliography}{5}

\setlength{\itemsep}{0.12em}
\setlength{\parskip}{0pt}

\bibitem{graphblas-api}
GraphBLAS Forum.
\textit{The GraphBLAS C API Specification. Version 2.1}, 2023.

\url{https://graphblas.org/docs/GraphBLAS_API_C_v2.1.0.pdf}

\bibitem{graphblas}
Davis, T. A.
\textit{SuiteSparse:GraphBLAS}.
Официальный репозиторий проекта.

\url{https://github.com/DrTimothyAldenDavis/GraphBLAS}

\bibitem{suitesparse}
Davis, T. A.
\textit{SuiteSparse}.
Официальный репозиторий проекта.

\url{https://github.com/DrTimothyAldenDavis/SuiteSparse}

\bibitem{ssmc}
Davis, T. A., Hu, Y.
\textit{The University of Florida Sparse Matrix Collection}.
ACM Transactions on Mathematical Software, 38(1), 2011.

\url{https://sparse.tamu.edu/}

\bibitem{cmake-ctest}
Kitware.
\textit{CTest Documentation}.

\url{https://cmake.org/cmake/help/latest/manual/ctest.1.html}

\end{thebibliography}

\end{document}
